\begin{chapintro}
統合した結果は、集めた試験に偏りがあれば偏ったままです。発表されなかった試験が抜け落ちていないかをファネルプロットと Egger 検定で調べ、解析の選び方で結論が変わらないかを感度分析で確かめます。最後に GRADE の考え方で、降圧薬 A の結果がどこまで確かかを判定し、結果の要約表にまとめます。固定効果モデルとランダム効果モデルの比較、累積メタ解析、1000 人あたりの人数を描く図は発展の節に置いてあり、初読では飛ばせます。
\end{chapintro}

\section{出版バイアスと小規模研究効果}
\label{sec:07-1}

有意な結果が出た試験は論文になりやすく、有意差の出なかった試験は投稿されなかったり、掲載が遅れたり、知名度の低い雑誌に載ったりします。その結果、文献データベースで見つかる試験は、行われた試験全体よりも効果が大きい側に偏ります。これを出版バイアスと呼びます。メタ解析は見つかった試験しか統合できないので、出版バイアスがあれば統合した効果も大きく出ます。

出版バイアスは、小さな試験で特に起こりやすい偏りです。大きな試験は費用と手間がかかっているので、結果にかかわらず発表されることが多い一方、小さな試験は効果が大きく出たときだけ発表される傾向があります。このため、小さな試験ほど効果が大きく出る傾向（小規模研究効果）が、出版バイアスの手がかりになります。ただし、小規模研究効果の原因は出版バイアスだけではありません。小さな試験は方法の質が低いことが多く（割り付けの隠蔽が不十分、など）、その分だけ効果が大きく出ることがあります。対象の重症度が違う、介入の強さが違う、といった本当の違いが原因のこともあります。小規模研究効果を見つけても、それがすなわち出版バイアスだとは言えない点に注意します。

\section{ファネルプロットを読む}
\label{sec:07-2}

小規模研究効果を目で見る図がファネルプロットです。横軸に各試験の効果（ここではリスク比、対数目盛）、縦軸に標準誤差をとって試験を点で置きます。縦軸は上ほど標準誤差（\kkref{sec:05-1-1}）が小さく、つまり大きな試験が上に来るように逆向きにとります。

\begin{rcode}
funnel(m, common = FALSE,
       xlab = "リスク比（対数目盛）")
\end{rcode}

\kkfig[.78\linewidth]{p1_funnel.pdf}{図 7.1　降圧薬 A の 10 試験のファネルプロット（架空のデータ）。縦の破線はランダム効果モデルの統合値、斜めの破線はその値から 95\% の範囲}

偏りがなければ、大きな試験は統合値の近くに集まり、小さな試験は左右に広く散らばるので、点は上がとがった漏斗（funnel）を逆さにした形に、左右対称に並びます。小さな試験のうち効果の小さい側（図の右下）が発表されずに抜けていると、左右が非対称になります。

図 7.1 では、下の方にある小さな試験（Endo 2015、Aoki 2009、Hayashi 2019）がそろって左側、つまり効果の大きい側にあり、右下はがら空きです。小さくて効果のない試験が抜けているのかもしれない、という読み方ができます。ただし、点が 10 個しかないと、偶然でもこの程度の非対称はよく起こります。目で見た印象は人によって違うので、数で確かめる方法が次の Egger 検定です。

\section{Egger 検定}
\label{sec:07-3}

Egger 検定は、ファネルプロットの非対称を数で見る検定です（Egger 1997）。試験が 10 本未満では当てになりません。

\begin{rcode}
metabias(m, method.bias = "Egger", k.min = 10)
\end{rcode}

\begin{rout}
Linear regression test of funnel plot asymmetry

Test result: t = -1.17, df = 8, p-value = 0.2741
Bias estimate: -1.2370 (SE = 1.0535)
\end{rout}

p = 0.27 で、非対称さは有意ではありませんでした。しかし、この結果から「出版バイアスはない」とは言えません。Egger 検定は試験数が少ないと検出力が低く、本当に偏りがあっても見逃しやすい検定です。目安として、試験数が 10 本未満ならファネルプロットの検定は行わないことが勧められており（Sterne 2011）、meta パッケージも既定では 10 本未満のときに計算しません（上の k.min = 10 はその下限です）。降圧薬 A の 10 試験は、ちょうどその下限にあたります。

\begin{hattenbox}
Egger 検定は、ファネルプロットの非対称さを回帰で調べます。各試験の効果を標準誤差で割った値（どれだけ 0 から離れているか）を、精度（標準誤差の逆数）に対して回帰し、切片が 0 からずれているかを検定します。小さな試験ほど効果が大きく出ていれば、切片が 0 からずれます。試験の大きさがどれも似ているときや、異質性が大きいときも、この検定は当てになりません。効果指標がオッズ比のときは、Egger 検定そのものが勧められていません。Cochrane ハンドブック（13.3.4.4）は、オッズ比では Harbord の検定や Peters の検定を使うよう勧めています（metabias の method.bias で選べます）。降圧薬 A はリスク比なので、Egger 検定のまま進めます。
\end{hattenbox}

\section{\hatten 固定効果とランダム効果の比較、選択モデル、trim and fill}
\label{sec:07-4}

Egger 検定のほかに、小規模研究効果の大きさを調べる方法として、固定効果モデルとランダム効果モデルの統合値を比べる方法があります。ランダム効果モデルは小さな試験に相対的に大きい重みを与えます。\kkref{sec:05-4} の数値では、Endo 2015 の重みは固定効果モデルの 1.3\% からランダム効果モデルの 2.9\% に増え、Garcia 2018 は 23.9\% から 16.5\% に減りました。そのため、小さな試験ほど効果が大きく出ているときは、ランダム効果モデルの統合値のほうが固定効果モデルより効果を大きく見せます。この比較が使えるのは、試験の間にばらつきがある（τ² > 0）ときだけです（Cochrane ハンドブック 13.3.4.6）。ばらつきがなければ、二つのモデルは同じ値になり、比べる意味がありません。

\begin{rcode}
round(exp(c(common = m$TE.common,
            random = m$TE.random)), 2)
\end{rcode}

\begin{rout}
common random
  0.75   0.73
\end{rout}

降圧薬 A では固定効果モデルが 0.75（0.67〜0.83）、ランダム効果モデルが 0.73（0.61〜0.86）で、差は小さく、向きもランダム効果のほうが効果を大きく見せる側ですが、結論を変える大きさではありません。この比較は、小規模研究効果の影響を大まかに確かめる感度分析として使います。このほか、Cochrane ハンドブックは、出版されやすさを統計モデルで表す選択モデルと、標準誤差などの試験の特徴で効果を回帰する方法を挙げています。いずれも専門家と相談して使う方法なので、この冊子では名前の紹介にとどめます。

論文では trim and fill という方法をよく見かけます。ファネルプロットが左右対称になるように、欠けていると推定される試験を反対側に補って統合し直す方法です。この冊子では使いません。非対称の原因がすべて出版バイアスで、欠けた試験がちょうど対称の位置にあるという強い仮定に立ち、異質性が大きいと性能が悪く、補った試験は実在しないので、補った後の値を補正後の効果として読むことはできないからです。

ファネルプロットと検定が調べているのは小規模研究効果であり、出版バイアスがあるかないかをそれだけで決めることはできません。Cochrane ハンドブック第 13 章（13.3）は、ROB-ME（risk of bias due to missing evidence、Page 2023）という手順を示しています。研究の中で報告されなかった結果と、研究の間で発表されなかった試験の両方について、結果が欠けた理由が結果の向きや有意性に左右されていないかを、構造的に判定する手順です。この判定は、この後の GRADE の出版バイアスの評価につながります。等高線つきファネルプロットは、非対称の原因を見分ける助けになるとして同じ節（13.3.4.3）で取り上げられていますが、この冊子では扱いません。

\section{未発表の試験を探す}
\label{sec:07-5}

出版バイアスへの対策として最も確かなのは、統計の工夫ではなく、未発表の試験を探すことです。ClinicalTrials.gov や jRCT などの試験登録には、結果が発表されていない試験も登録されています。登録されているのに結果が見つからない試験が多ければ、それ自体が出版バイアスの証拠になります。\kkref{sec:03-1} で試験登録を検索先に入れたのはこのためです。

\section{感度分析}
\label{sec:07-6}

感度分析は、解析の途中で行った選択を変えても結論が変わらないかを確かめる作業です。第 2 部で検査の「感度」が出てきますが、それとは別の言葉です。検査の感度は病気の人を陽性と判定する割合で、メタ解析の感度分析は結果が選択にどれだけ敏感かを調べることを指します。

降圧薬 A では、計画書で次の感度分析を決めておいたとします。

\begin{description}
\item[バイアスのリスクが高い試験を除く] RoB 2 で全体の判定が高かった Endo 2015 を除くと、リスク比は 0.74（0.62〜0.88）で、ほとんど変わりませんでした。
\item[1 試験ずつ除く] 10 試験から 1 本ずつ除いて 10 回統合し直すと、リスク比は 0.70（Chen 2012 または Dubois 2014 を除いたとき）から 0.77（Jensen 2022 を除いたとき）の範囲に収まりました。特定の 1 試験が結論を左右してはいません。
\item[モデルと方法を変える] 固定効果モデルでは 0.75（0.67〜0.83）、Hartung-Knapp 法（\kkref{sec:05-7}）では 0.73（0.59〜0.89）で、どれも 1 を含みませんでした。
\item[出版バイアスを考慮する] 固定効果モデル（0.75）とランダム効果モデル（0.73）の統合値を比べました。ランダム効果モデルでは小さな試験の重みが増えます（\kkref{sec:05-4} で、Endo 2015 は 1.3\% から 2.9\%）が、統合値の差は小さく、結論は動きませんでした。この比較は発展の \kkref{sec:07-4} で説明しています。
\end{description}

\begin{rcode}
update(m, subset = rob != "高")   # RoB 高を除く
metainf(m, pooled = "random")    # 1 試験ずつ除く
\end{rcode}

subset には残す試験の条件を書きます。!= は「等しくない」という意味です。metainf の結果は forest に渡すと、1 試験ずつ除いた結果のフォレストプロットになります。

\begin{rcode}
mi <- metainf(m, pooled = "random")
forest(mi)
\end{rcode}

\kkfig[.9\linewidth]{x1_loo.pdf}{図 7.2　1 試験ずつ除いたときの統合したリスク比（ランダム効果モデル）。各行は、その試験を除いた 9 試験の結果}

図 7.2 の各行は、左の列の試験を除いた 9 試験で統合し直した結果です。1 行が四角と横線で示され、最下段が 10 試験すべての結果（0.73、0.61〜0.86）です。四角の位置は 0.70（Chen 2012 または Dubois 2014 を除いたとき）から 0.77（Jensen 2022 を除いたとき）の間にあり、どの行でも区間は 1 をまたいでいません。Jensen 2022 を除くと、点推定は最も 1 に近づきます。効果の大きい試験を除いたためです。Chen 2012 と Dubois 2014 を除くと逆に 0.70 まで下がり、効果が小さい（点推定が 1 に近いか 1 より大きい）試験を除いたことになります。見るのは、全体の結果から大きく離れる行があるか、区間が 1 をまたぐ行があるかで、この例ではどちらもありません。

どの感度分析でもリスク比は 0.70〜0.77 の範囲にあり、信頼区間が 1 を含むことはありませんでした。効果の向きと大きさについての結論は、解析の選び方に対して頑健だと言えます。もし感度分析で結論が変わったら、その結果を隠さずに報告し、どちらの結果をどう解釈するかを考察で論じます。

\section{\hatten 累積メタ解析}
\label{sec:07-7}

累積メタ解析（\kkref{sec:01-3-1} の図 1.1）は、metacum で、発表年の順に足した結果を出せます。

\begin{rcode}
year <- as.numeric(sub(".* ", "", bp$study))  # 試験名から年
mc <- metacum(m, sortvar = year, pooled = "random")
mc
forest(mc)
\end{rcode}

\begin{rout}
                           RR            95% CI p-value
Adding Aoki 2009 (k=1)    0.4118 [0.1759; 0.9640]  0.0409
Adding Berger 2011 (k=2)  0.6290 [0.3774; 1.0485]  0.0754
Adding Chen 2012 (k=3)    0.7756 [0.4829; 1.2456]  0.2931
Adding Dubois 2014 (k=4)  0.8704 [0.7086; 1.0692]  0.1860
Adding Endo 2015 (k=5)    0.8343 [0.6613; 1.0525]  0.1265
Adding Fischer 2016 (k=6) 0.8625 [0.7363; 1.0103]  0.0669
Adding Garcia 2018 (k=7)  0.7815 [0.6455; 0.9462]  0.0115
Adding Hayashi 2019 (k=8) 0.7529 [0.6184; 0.9168]  0.0047
Adding Ivanov 2021 (k=9)  0.7690 [0.6590; 0.8975]  0.0009
Adding Jensen 2022 (k=10) 0.7281 [0.6137; 0.8637]  0.0003
\end{rout}

sortvar に年を渡すと、その順に試験を足します。k = 2 から 6 の行は区間の上限が 1 を超えており、k = 7 から超えなくなります。

\section{GRADE で確実性を判定する}
\label{sec:07-8}

ここまでの結果を、診療の判断に使えるかどうかの言葉に直すのが GRADE です（Guyatt 2008）。GRADE は、あるアウトカムについての効果の推定値がどれくらい確かかを、高、中、低、非常に低の 4 段階で判定する方法で、多くの診療ガイドラインで使われています。

RCT だけを集めたメタ解析は、確実性が高から出発します。そこから、次の五つの理由があれば 1 段階か 2 段階下げます。

\begin{description}
\item[バイアスのリスク] 統合した試験の方法に問題があるか。RoB 2 の判定を、重みの大きい試験を中心に見ます。
\item[非一貫性] 試験の結果が食い違っているか。I²（\kkref{sec:06-1}）の数字だけで決めず、効果の向きと大きさのずれ、信頼区間の重なり、予測区間が判断の閾値（下の不精確さで説明します）をまたぐかを見ます（Core GRADE 3）。信頼区間と予測区間の違いは \kkref{sec:06-3} で説明しました。
\item[非直接性] 試験の対象、介入、比較、アウトカムが、知りたい問いとずれていないか。
\item[不精確さ] 信頼区間が広すぎて、判断が変わる範囲にまたがっていないか。判断の閾値は、結論が変わる境目の効果の大きさです。効果なし（リスク比 1）を置くことも、最小重要差（患者や臨床医が意味があると考える最小の差）を置くこともあります。閾値を先に決めておき、信頼区間がその閾値をまたぐかで判断します（Core GRADE 2）。事象の数が少なすぎないかは補助として見ます。
\item[出版バイアス] 発表されていない試験が抜けている疑いが強いか。
\end{description}

観察研究のメタ解析は低から出発し、効果が非常に大きい、量反応関係がある、といった理由があれば上げることができます（ROBINS-I で評価する場合は高から出発し、バイアスのリスクの評価で下げる扱いもあります。Cochrane ハンドブック 14.2.1）。

降圧薬 A の脳卒中について判定すると、次のようになります。

\begin{description}
\item[バイアスのリスク] 重みの大きい Garcia 2018、Dubois 2014、Ivanov 2021 はいずれも低で、高の Endo 2015 を除いても結果は変わらないので、下げません。
\item[非一貫性] 点推定は Aoki 2009 の 0.41 から、Chen 2012 の 1 を超える 1.13 までずれています。10 試験のうち 9 試験は 1 より小さいものの、Chen 2012 だけが逆向きです。予測区間は 0.46〜1.15 で、判断の閾値である 1（効果なし）をまたぎます。I² は 51\% でした。サブグループ解析でも説明できなかったので、1 段階下げます。
\item[非直接性] 対象は高血圧の成人、比較はプラセボ、アウトカムは脳卒中で、問いと一致しているので下げません。実際の診療では他の降圧薬と比べたいことが多いので、プラセボとの比較であることは結果の要約表に書き添えます。
\item[不精確さ] 判断の閾値を効果なし（リスク比 1）に置くと、信頼区間 0.61〜0.86 はその片側に収まります。対照群のリスクが 5\% なら、1000 人あたりの差は 1000 × 0.05 × (1 − リスク比) で求められ、区間の端のリスク比 0.864 と 0.614 を入れると 7〜19 人の減少で（表 7.1）、区間のどこをとっても減る側です。最小重要差を置くなら、1000 人あたり何人の減少から臨床的に意味があるとするかを先に決め、7〜19 人がその値の片側に収まるかを見ます。事象も 1,304 件あるので、下げません。
\item[出版バイアス] 小さな試験が効果の大きい側に偏っていますが、Egger 検定は有意でなく、固定効果とランダム効果の統合値も近い値（0.75 と 0.73、発展の \kkref{sec:07-4}）でした。疑いはあるものの下げるほどではないと判断し、その旨を記録します。
\end{description}

判定の結果、確実性は高から 1 段階下げて中となります。判定は機械的な点数の計算ではなく、レビューの著者の判断です。だからこそ、どの理由でいくつ下げたかを必ず書き残し、読者が判断の妥当性を確かめられるようにします。

\section{結果の要約表}
\label{sec:07-9}

GRADE では、アウトカムごとの結果を Summary of findings 表（結果の要約表）にまとめます。相対的な効果（リスク比）だけでなく、絶対的な効果（1000 人あたり何人減るか）を示すのが特徴です。リスク比は集団によらず比較的一定でも、もともとの発症割合が違えば、減る人数は大きく変わるからです。

\begin{kkcenter}
\kkfit{\begin{tabular}{lccc}
\toprule
対照群のリスク & プラセボ & 降圧薬 A & 差（95\% 信頼区間） \\
\midrule
5\%（一般の成人） & 1000 人中 50 人 & 1000 人中 36 人 & 14 人少ない（7〜19 人少ない） \\
10\%（高齢者） & 1000 人中 100 人 & 1000 人中 73 人 & 27 人少ない（14〜39 人少ない） \\
\bottomrule
\end{tabular}}
\kkcap{表 7.1　降圧薬 A と脳卒中の結果の要約（架空のデータ）。リスク比 0.73（0.61〜0.86）、10 試験、21,420 人、確実性は中（非一貫性で 1 段階下げた）}
\end{kkcenter}

対照群のリスクは、統合した試験の対照群の発症割合（およそ 5〜13\%）を目安に、低い場合と高い場合の 2 通りを置きました。降圧薬 A 群の人数は、対照群のリスクにリスク比 0.73 をかけて求めます。1000 人あたり 50 人なら 36 人に、100 人なら 73 人になります。同じリスク比でも、もともと脳卒中を起こしやすい高齢者では、防げる人数がおよそ 2 倍になることが表から読み取れます。対照群のリスクを連続的に動かした図は、発展の \kkref{sec:07-9-1} にあります。

\subsection{\hatten 対照群のリスクを動かした図}
\label{sec:07-9-1}

対照群のリスクを連続的に動かすと、防げる人数がどう変わるかを一枚の図にできます。リスク比 0.73 とその信頼区間の端を使い、対照群のリスク p について 1000 × p × (1 − リスク比) を計算して線にします。

\begin{rcode}
rr <- exp(c(m$TE.random, m$lower.random, m$upper.random))
p0 <- seq(0.01, 0.15, by = 0.001)    # 対照群の 5 年発症割合
plot(100 * p0, 1000 * p0 * (1 - rr[1]), type = "l",
     ylim = c(0, 65), xlab = "プラセボ群の発症割合（%）",
     ylab = "1000 人あたり防げる脳卒中（人）")
lines(100 * p0, 1000 * p0 * (1 - rr[2]), lty = 2)  # 下限
lines(100 * p0, 1000 * p0 * (1 - rr[3]), lty = 2)  # 上限
\end{rcode}

\kkfig[.8\linewidth]{x1_absolute.pdf}{図 7.3　プラセボ群の 5 年の発症割合と、1000 人あたり防げる脳卒中の人数。帯は 95\% 信頼区間。赤い点は表 7.1 の 5\% と 10\%}

横軸はプラセボ群での脳卒中の発症割合、縦軸は降圧薬 A によって 1000 人あたり防げる脳卒中の人数で、中央の線が点推定、帯が 95\% 信頼区間です。赤い点が表 7.1 の 2 つの行にあたり、5\% なら 14 人、10\% なら 27 人です。線は原点から右上がりの直線で、リスク比が一定なら、もともとの発症割合が高いほど防げる人数が増えます。帯は右へ行くほど広がり、10\% で 14〜39 人です。

読むときの注意は、この線が降圧薬 A のリスク比 0.73 を、どの対照群のリスクにもそのまま当てはめた計算だということです。\kkref{sec:06-3} で見たとおり、効果は集団によって違いうるので、予測区間が 1 をまたぐことを忘れて線を読むと、確かさを過大に見積もります。また、1000 人あたりの人数は期間を決めて初めて意味をもつので、横軸の発症割合が何年間の割合なのかを、図と一緒に書いておく必要があります。

\section{この章のまとめ}
\label{sec:07-10}

ファネルプロットでは小さな試験が効果の大きい側に寄っていましたが、Egger 検定は試験数 10 で検出力が低く、p = 0.27 でした。固定効果モデル（0.75）とランダム効果モデル（0.73）の統合値の差は小さく、結論は変わりませんでした。バイアスのリスク、1 試験ずつの除外、モデルの選択を変えた感度分析でも、リスク比は 0.70〜0.77 に収まりました。GRADE の判定は非一貫性で 1 段階下げて中とし、結果の要約表に 1000 人あたりの差で示しました。小規模研究効果は出版バイアス以外の原因でも起こり、統計の方法で出版バイアスを取り除くことはできないので、試験登録の検索が最も確かな対策です。次の章では、ここまでの作業を論文として報告する方法を説明します。
